\begin{definition}{Система $M^{[X]}|M|1$}
Система массового обслуживания, на вход которой поступает пуассоновский поток событий с интенсивностью $\lambda$. В данной модели каждое событие инициирует поступление не одиночной заявки, а группы (пакета). 

Размер группы является дискретной случайной величиной $X$. Предполагается, что размеры последовательно поступающих групп представляют собой независимые в совокупности и одинаково распределенные случайные величины. Закон распределения размера группы задается вероятностями:
$$ \P(X=n) = c_n, \quad n \in \N, $$
где $\dsum_{n=1}^{\infty} c_n = 1$. Интенсивность поступления групп, состоящих строго из $n$ заявок, составляет $\lambda_n = \lambda c_n$. Время обслуживания отдельной заявки распределено по экспоненциальному закону с параметром $\mu$.
\end{definition}

Вследствие группового поступления заявок граф переходов марковской цепи содержит дуги, соответствующие скачкам процесса на произвольное число состояний вправо.

\begin{figure}[h!]
    \centering
    \begin{tikzpicture}[
        scale=1.5,
        >={Stealth[length=3mm]},
        state/.style={circle, draw, minimum size=0.8cm, thick},
        every node/.style={font=\small}
    ]
        % Состояния
        \node[state] (s0) at (0,0) {0};
        \node[state] (s1) at (2,0) {1};
        \node[state] (s2) at (4,0) {2};
        \node[state] (s3) at (6,0) {3};
        \node (dots) at (7.5,0) {\Large $\dots$};

        % Переходы \lambda_1
        \draw[->, thick] (s0) to [bend left=30] node[midway, fill=Snow2, inner sep=1.5pt] {$\lambda_1$} (s1);
        \draw[->, thick] (s1) to [bend left=30] node[midway, fill=Snow2, inner sep=1.5pt] {$\lambda_1$} (s2);
        \draw[->, thick] (s2) to [bend left=30] node[midway, fill=Snow2, inner sep=1.5pt] {$\lambda_1$} (s3);

        % Переходы \lambda_2
        \draw[->, thick] (s0) to [bend left=50] node[midway, fill=Snow2, inner sep=1.5pt] {$\lambda_2$} (s2);
        \draw[->, thick] (s1) to [bend left=50] node[midway, fill=Snow2, inner sep=1.5pt] {$\lambda_2$} (s3);

        % Обслуживание \mu
        \draw[->, thick] (s1) to [bend left=30] node[midway, below] {$\mu$} (s0);
        \draw[->, thick] (s2) to [bend left=30] node[midway, below] {$\mu$} (s1);
        \draw[->, thick] (s3) to [bend left=30] node[midway, below] {$\mu$} (s2);
    \end{tikzpicture}
    \caption{Граф интенсивностей переходов для системы с групповым поступлением заявок. Проиллюстрированы приходы групп размера 1 и 2.}
\end{figure}

\begin{statement}
Производящая функция стационарного распределения вероятностей состояний системы $P(z) = \dsum_{n=0}^{\infty} p_n z^n$ имеет вид:
$$ P(z) = \frac{p_0}{1 - \frac{\lambda}{\mu} z \overline{C}(z)}, $$
где $\overline{C}(z) = \frac{1 - C(z)}{1 - z}$, $C(z)$ — производящая функция распределения размера группы $X$.
Вероятность простоя системы равна $p_0 = 1 - \rho$, где коэффициент загрузки $\rho = \frac{\lambda}{\mu}\E[X] < 1$.

Стационарное распределение полностью определяется данной функцией: вероятности состояний $\{p_n\}_{n=0}^{\infty}$ вычисляются как коэффициенты разложения $P(z)$ в ряд Тейлора:
$$ p_n = \frac{1}{n!} \left. \frac{d^n P(z)}{dz^n} \right|_{z=0}. $$
\end{statement}

\begin{sproof}
Составим систему уравнений глобального баланса, приравнивая интенсивности выхода из состояния $n$ и входа в него:
\begin{itemize}
    \item Для состояния $n = 0$:
    $$ \lambda p_0 = \mu p_1. $$
    \item Для состояний $n \ge 1$:
    $$ (\lambda + \mu) p_n = \mu p_{n+1} + \dsum_{k=1}^{n} p_{n-k} \lambda_k. $$
\end{itemize}
Подставляя $\lambda_k = \lambda c_k$, перепишем второе уравнение:
$$ (\lambda + \mu) p_n = \mu p_{n+1} + \lambda \dsum_{k=1}^{n} p_{n-k} c_k. $$
Введем производящие функции для распределения размера группы $C(z) = \dsum_{n=1}^{\infty} c_n z^n$ и вероятностей состояний $P(z) = \dsum_{n=0}^{\infty} p_n z^n$. Умножим уравнение для $n \ge 1$ на $z^n$ и просуммируем по $n$ от $1$ до $\infty$:
$$ (\lambda + \mu) \dsum_{n=1}^{\infty} p_n z^n = \mu \dsum_{n=1}^{\infty} p_{n+1} z^n + \lambda \dsum_{n=1}^{\infty} \dsum_{k=1}^{n} p_{n-k} c_k z^n. $$
Выразим одинарные суммы через $P(z)$:
$$ \dsum_{n=1}^{\infty} p_n z^n = P(z) - p_0, \quad \dsum_{n=1}^{\infty} p_{n+1} z^n = \frac{1}{z} \dsum_{n=2}^{\infty} p_n z^n = \frac{1}{z} (P(z) - p_0 - p_1 z). $$
Для преобразования двойной суммы изменим порядок суммирования. Исходная область индексов представляет собой дискретный клин $1 \le k \le n < \infty$.

\begin{center}
\begin{tikzpicture}[scale=1.2, >=Stealth, font=\small]
    \draw[->, thick] (-0.2,0) -- (4,0) node[right] {$n$};
    \draw[->, thick] (0,-0.2) -- (0,4) node[above] {$k$};
    
    \draw[thick, DeepSkyBlue4] (0,0) -- (3.5,3.5) node[above right] {$k=n$};
    
    \fill[pattern=north west lines, pattern color=DeepSkyBlue4, opacity=0.5] (1,1) -- (3.5,1) -- (3.5,3.5) -- cycle;
    \draw[dashed] (1,0) -- (1,1);
    \node[below] at (1,0) {$1$};
    \draw[dashed] (0,1) -- (1,1);
    \node[left] at (0,1) {$1$};
    
    \node[fill=Snow2, inner sep=2pt] at (2.5, 1.5) {$k \le n$};
\end{tikzpicture}
\end{center}

Меняя пределы, получаем произведение двух рядов:
$$ \dsum_{n=1}^{\infty} \dsum_{k=1}^{n} p_{n-k} c_k z^n = \dsum_{k=1}^{\infty} \dsum_{n=k}^{\infty} c_k p_{n-k} z^n = \dsum_{k=1}^{\infty} c_k z^k \dsum_{n=k}^{\infty} p_{n-k} z^{n-k} = C(z) P(z). $$
Подставим полученные выражения в основное уравнение:
$$ (\lambda + \mu)(P(z) - p_0) = \frac{\mu}{z}(P(z) - p_0 - p_1 z) + \lambda C(z) P(z). $$
Так как $\frac{\mu}{z} p_1 z = \mu p_1$, и из уравнения для $n=0$ известно соотношение $\mu p_1 = \lambda p_0$, получаем:
$$ \lambda P(z) - \lambda p_0 + \mu(P(z) - p_0) = \frac{\mu}{z}(P(z) - p_0) - \lambda p_0 + \lambda C(z) P(z). $$
После сокращения слагаемых $-\lambda p_0$ группируем члены, содержащие $P(z)$, в левой части, а $p_0$ — в правой:
$$ \lambda P(z) - \lambda C(z) P(z) = \frac{\mu}{z}(P(z) - p_0) - \mu(P(z) - p_0). $$
Вынося общие множители, находим:
$$ \lambda P(z) (1 - C(z)) = \mu (P(z) - p_0) \left( \frac{1 - z}{z} \right). $$
Умножая обе части на $z$, выражаем $P(z)$:
$$ \lambda z P(z) (1 - C(z)) = \mu (1 - z) P(z) - \mu (1 - z) p_0, $$
$$ P(z) \left[ \mu(1 - z) - \lambda z (1 - C(z)) \right] = \mu(1 - z) p_0, $$
$$ P(z) = \frac{\mu(1 - z) p_0}{\mu(1 - z) - \lambda z (1 - C(z))}. $$
Разделив числитель и знаменатель на $\mu(1 - z)$ и введя функцию $\overline{C}(z) = \frac{1 - C(z)}{1 - z}$, приходим к требуемому соотношению для производящей функции.

Найдем вероятность простоя $p_0$, используя условие нормировки $P(1) = 1$. Перейдем к пределу при $z \to 1$. Значение $\overline{C}(1)$ вычисляется по правилу Лопиталя:
$$ \lim_{z \to 1} \overline{C}(z) = \lim_{z \to 1} \frac{1 - C(z)}{1 - z} = \lim_{z \to 1} \frac{-C'(z)}{-1} = C'(1). $$
Поскольку $C'(1) = \dsum_{n=1}^{\infty} n c_n = \E[X]$, знаменатель функции $P(z)$ в точке $z=1$ принимает вид $1 - \frac{\lambda}{\mu} \E[X]$. Из условия нормировки следует:
$$ 1 = \frac{p_0}{1 - \frac{\lambda}{\mu} \E[X]} \implies p_0 = 1 - \rho, $$
где $\rho = \frac{\lambda}{\mu} \E[X]$. Для существования стационарного режима необходимо выполнение условия $p_0 > 0$, что эквивалентно $\rho < 1$.
\end{sproof}

\begin{theorem}{Макроскопические характеристики $M^{[X]}|M|1$}
При условии $\rho < 1$ стационарный режим существует. Основные метрики эффективности системы определяются соотношениями:
\begin{itemize}
    \item Среднее число заявок в системе: $\displaystyle L = \frac{\rho + \frac{\lambda}{\mu}\E[X^2]}{2(1 - \rho)}$.
    \item Среднее время пребывания заявки: $\displaystyle W = \frac{L}{\lambda \E[X]}$.
    \item Среднее время ожидания: $\displaystyle W_q = W - \frac{1}{\mu}$.
    \item Среднее число заявок в очереди: $\displaystyle L_q = L - \rho$.
\end{itemize}
\end{theorem}

\begin{tproof}
Среднее число заявок $L$ вычисляется как значение первой производной производящей функции в точке $1$: $L = P'(1)$. Представим $P(z)$ в виде $P(z) = p_0 \left( 1 - \frac{\lambda}{\mu} z \overline{C}(z) \right)^{-1}$ и продифференцируем:
$$ P'(z) = p_0 \left( 1 - \frac{\lambda}{\mu} z \overline{C}(z) \right)^{-2} \cdot \frac{\lambda}{\mu} \left( \overline{C}(z) + z \overline{C}'(z) \right). $$
Для нахождения предела $\overline{C}'(1)$ продифференцируем дробь $\overline{C}(z)$ по правилу дифференцирования частного:
$$ \overline{C}'(z) = \frac{-C'(z)(1-z) - (1-C(z))(-1)}{(1-z)^2} = \frac{-C'(z)(1-z) + 1 - C(z)}{(1-z)^2}. $$
При $z \to 1$ возникает неопределенность $\left[\frac{0}{0}\right]$. Применяем правило Лопиталя:
$$ \lim_{z \to 1} \overline{C}'(z) = \lim_{z \to 1} \frac{-C''(z)(1-z) - C'(z)(-1) - C'(z)}{-2(1-z)} = \lim_{z \to 1} \frac{-C''(z)(1-z)}{-2(1-z)} = \frac{1}{2} C''(1). $$
По определению второй факториальный момент равен $C''(1) = \E[X(X-1)] = \E[X^2] - \E[X]$. 
Подставляя найденные пределы в выражение для $P'(1)$, получаем:
$$ L = P'(1) = \frac{p_0}{p_0^2} \cdot \frac{\lambda}{\mu} \left( \E[X] + \frac{1}{2}(\E[X^2] - \E[X]) \right) = \frac{1}{1-\rho} \cdot \frac{\lambda}{2\mu} \left( \E[X] + \E[X^2] \right). $$
Учитывая, что $\frac{\lambda}{2\mu}\E[X] = \frac{\rho}{2}$, итоговая формула длины очереди принимает вид:
$$ L = \frac{\rho + \frac{\lambda}{\mu}\E[X^2]}{2(1-\rho)}. $$

Остальные метрики выводятся на основе закона Литтла. Поскольку группы заявок поступают с интенсивностью $\lambda$, а средний размер одной группы равен $\E[X]$, эффективная интенсивность поступления индивидуальных заявок (общее число единичных заявок, приходящих в систему за единицу времени) определяется произведением: 
$$ \lambda_{\text{eff}} = \lambda \E[X]. $$
Отсюда среднее время пребывания $W = \frac{L}{\lambda_{\text{eff}}}$. Средняя длина очереди $L_q = \lambda_{\text{eff}} W_q = \lambda_{\text{eff}} \left(W - \frac{1}{\mu}\right) = L - \frac{\lambda_{\text{eff}}}{\mu} = L - \rho$.
\end{tproof}

\begin{example}{Фиксированный размер группы}
Рассмотрим систему $M^{[X]}|M|1$, в которой заявки поступают строго пакетами фиксированной длины $k$ ($X = k = \text{const}$). Необходимо вычислить среднее число заявок в системе $L$.
\end{example}

\begin{eproof}
Для детерминированной величины $X = k$ начальные моменты равны $\E[X] = k$ и $\E[X^2] = k^2$. Коэффициент загрузки составляет $\rho = \frac{\lambda k}{\mu}$.
Подставляя моменты в полученную формулу для макроскопических характеристик, находим:
$$ L = \frac{\rho + \frac{\lambda}{\mu} k^2}{2(1 - \rho)} = \frac{\rho + \rho k}{2(1 - \rho)} = \frac{k+1}{2} \cdot \frac{\rho}{1 - \rho}. $$
\end{eproof}

\begin{example}{Геометрическое распределение размера группы}
Пусть размер группы имеет геометрическое распределение с параметром $\alpha \in (0, 1)$: $\P(X = n) = (1 - \alpha) \alpha^{n-1}$ для $n \in \N$. Требуется определить среднее число заявок в системе $L$.
\end{example}

\begin{eproof}
Составим производящую функцию $C(z)$:
$$ C(z) = \dsum_{n=1}^{\infty} (1 - \alpha) \alpha^{n-1} z^n = \frac{z(1 - \alpha)}{1 - \alpha z}. $$
Вычислим первый момент $\E[X]$:
$$ \E[X] = \lim_{z \to 1} \frac{1 - C(z)}{1 - z} = \lim_{z \to 1} \frac{1 - \frac{z(1-\alpha)}{1-\alpha z}}{1 - z} = \lim_{z \to 1} \frac{1 - \alpha z - z + \alpha z}{(1 - \alpha z)(1 - z)} = \lim_{z \to 1} \frac{1 - z}{(1 - \alpha z)(1 - z)} = \frac{1}{1 - \alpha}. $$
Второй начальный момент для смещенного геометрического распределения равен $\E[X^2] = \frac{1 + \alpha}{(1 - \alpha)^2}$. Загрузка системы составит $\rho = \frac{\lambda}{\mu (1 - \alpha)}$.
Подставим моменты в общую формулу:
$$ L = \frac{\rho + \frac{\lambda}{\mu} \frac{1 + \alpha}{(1 - \alpha)^2}}{2(1 - \rho)}. $$
Так как $\frac{\lambda}{\mu} = \rho (1 - \alpha)$, преобразуем числитель:
$$ L = \frac{\rho + \rho (1 - \alpha) \frac{1 + \alpha}{(1 - \alpha)^2}}{2(1 - \rho)} = \frac{\rho \left(1 + \frac{1 + \alpha}{1 - \alpha}\right)}{2(1 - \rho)}. $$
После приведения скобки к общему знаменателю $\left( \frac{1 - \alpha + 1 + \alpha}{1 - \alpha} = \frac{2}{1 - \alpha} \right)$, получаем итоговый результат:
$$ L = \frac{\rho \cdot \frac{2}{1 - \alpha}}{2(1 - \rho)} = \frac{1}{1 - \alpha} \cdot \frac{\rho}{1 - \rho}. $$
\end{eproof}

\begin{note}
В обоих рассмотренных примерах среднее число заявок в системе $L$ принимает вид произведения средней длины очереди классической системы $M|M|1$ (равной $\frac{\rho}{1-\rho}$) на константный множитель. Общую формулу для $L$ можно переписать в виде:
$$ L = \frac{\rho}{1 - \rho} \cdot \frac{\E[X^2] + \E[X]}{2 \E[X]}. $$
Данный мультипликативный фактор зависит исключительно от первых двух моментов распределения размера группы. Это отражает следующее свойство (аналогичное формуле Поллачека-Хинчина): при фиксированной общей загрузке $\rho$ дисперсия размера пакета напрямую увеличивает среднюю длину очереди. Чем больше разброс размера группы по сравнению с одиночным поступлением, тем на большую величину масштабируется базовая метрика системы $M|M|1$.
\end{note}